\Chapter{The Beckmann Theorem}\label{cap:beckmann}

\section{The Beckmann problem}

\subsection{Formulation}

\datebox{2026}

\begin{definition}[Beckmann problem]\label{def:beckmann}
The Beckmann problem is
\[
(\mathcal{B}) \quad \min_{f \in \FeasSet}
Z(f) \defeq \sum_{a \in A} T_a(x_a(f))
= \sum_{a \in A}\int_0^{x_a(f)}t_a(u)\,du.
\]
\end{definition}

\begin{proposition}\label{prop:beckmann_convex}
The Beckmann problem is convex. If each $t_a$ is strictly increasing, then $Z$ is strictly convex in edge flows.
\end{proposition}

\begin{remark}
The Beckmann objective is not the total system cost $\sum_a x_a t_a(x_a)$. This distinction separates user equilibrium from social optimum.
\end{remark}

\subsection{Main theorem}

\begin{theorem}[Beckmann, McGuire, and Winsten, 1956]\label{thm:beckmann}
For continuous, positive, and strictly increasing edge costs, $f^* \in \FeasSet$ is a Wardrop equilibrium if and only if it solves the Beckmann problem.
\end{theorem}

\begin{proof}
The proof follows from convexity and the KKT conditions. The gradient of $Z$ with respect to a route flow is the route cost:
\[
\frac{\partial Z}{\partial f_p} = \sum_a t_a(x_a)\delta_{ap} = C_p(x).
\]
The KKT conditions are therefore exactly the Wardrop complementarity conditions.
\end{proof}
